% GENERATED FILE. Edit canonical lessons, then run npm run generate:books.
% canonical-source-hash: fnv1a64:a4b0958d951be877
% canonical-lessons: RU-C268-bolshe-chem, RU-C268-uzhe-let, RU-C268-review

\chapter{Better Than, and Years Ago}
\label{ch:ru-better-than-and-years-ago}
\hlchaptermodality{268}

\begin{chapteropening}
\textbf{By the end of this chapter:} \emph{I can compare two things with -\ru{ее}, \ru{больше}, \ru{лучше} and \ru{чем}, and say how long and how long ago with \ru{уже} ... \ru{года} and ... \ru{назад}.}

Complete RU-C268-review: compare two ways to travel, then say how long you have done something and when you did it.
\end{chapteropening}

\section[\texorpdfstring{\ru{быстрее}, \ru{чем} … … (bystréye, chem …)}{быстрее, чем … … (bystréye, chem …)}]{\ru{быстрее}, \ru{чем}; \ru{лучше}, \ru{хуже} — comparing two things}

\label{lesson:RU-C268-bolshe-chem}



\begin{quote}
Say \textquotedblleft{}fast\textquotedblright{}, \textquotedblleft{}expensive\textquotedblright{} and \textquotedblleft{}cheap\textquotedblright{}. (\emph{Býstryy. Dorogóy. Deshyóvyy.}) Then \textquotedblleft{}you'd better take a taxi\textquotedblright{}. (\emph{Tebé lúchshe vzyat taksí.}) That \emph{\ru{лучше}} was already a comparison: \textquotedblleft{}better\textquotedblright{}.
\end{quote}

\begin{grammarlens}[title={Grammar Lens: -\ru{ее}, and \ru{чем}}]
Most adjectives compare by swapping their ending for \textbf{-\ru{ее}}: \emph{\ru{быстрый}}, \textbf{\ru{быстрее}} (\emph{bystréye}); \emph{\ru{тёплый}}, \textbf{\ru{теплее}} (\emph{tepléye}); \emph{\ru{интересный}}, \textbf{\ru{интереснее}} (\emph{interésneye}). The form does not change for gender or number.

\textquotedblleft{}Than\textquotedblright{} is \textbf{\ru{чем}}, with a comma in front:

\begin{quote}
\textbf{\ru{Поезд} \ru{быстрее}, \ru{чем} \ru{автобус}.} — \emph{póyezd bystréye, chem avtóbus} — The train is faster than the bus.
\end{quote}

In a short adjective the stress usually moves onto \emph{-\ru{ее}}: \emph{bystréye}. In a longer one it usually stays where it was: \emph{interésneye}.
\end{grammarlens}

\begin{grammarlens}[title={Grammar Lens: the ones that change shape}]
A few common ones are learned whole: \emph{\ru{большой}}, \textbf{\ru{больше}} (\emph{bólshe}), bigger, more; \emph{\ru{маленький}}, \textbf{\ru{меньше}} (\emph{ménshe}), smaller, less; \emph{\ru{хороший}}, \textbf{\ru{лучше}} (\emph{lúchshe}), better; \emph{\ru{плохой}}, \textbf{\ru{хуже}} (\emph{khúzhe}), worse; \emph{\ru{дорогой}}, \textbf{\ru{дороже}} (\emph{dorózhe}); \emph{\ru{дешёвый}}, \textbf{\ru{дешевле}} (\emph{deshévle}).

\begin{quote}
\textbf{\ru{Самолёт} \ru{быстрее}, \ru{но} \ru{поезд} \ru{дешевле}.} — \emph{samolyót bystréye, no póyezd deshévle} — The plane is faster, but the train is cheaper. \textbf{\ru{По}-\ru{моему}, \ru{поезд} \ru{лучше}.} — \emph{po-móyemu, póyezd lúchshe} — In my view the train is better.
\end{quote}
\end{grammarlens}

\begin{cousinweb}
\textbf{\ru{чем}} is a form of \textbf{\ru{что}}, the \textquotedblleft{}by what\textquotedblright{} form, on the question-root PIE \emph{*k\textsuperscript{w}o-} behind \emph{\ru{что}}, \emph{\ru{кто}} and English \emph{what}. EtymDB 2.1 lists \emph{\ru{чем}} without taking it apart, so the link to \emph{\ru{что}} is the plain reading of the forms.
\end{cousinweb}

\subsection*{Your turn}
\begin{itemize}
\raggedright
  \item \emph{Say it:} \emph{Póyezd bystréye, chem avtóbus.}
  \item \emph{Say it:} \emph{Samolyót bystréye, no póyezd deshévle.}
  \item \emph{Say it:} \emph{bólshe, ménshe, lúchshe, khúzhe}
  \item \emph{Answer:} \emph{Shto lúchshe: póyezd íli samolyót?} \textquotedblleft{}Which is better?\textquotedblright{} Give one reason with \emph{potomú shto}
\end{itemize}

\begin{tcolorbox}[breakable,colback=teal!4,colframe=teal!35!black,title={Before you move on}]
\textquotedblleft{}Warmer\textquotedblright{}? (\emph{\textbf{\ru{Теплее}.}}) \textquotedblleft{}Than\textquotedblright{}? (\emph{\textbf{\ru{Чем}}}, with a comma.) \textquotedblleft{}Worse\textquotedblright{}? (\emph{\textbf{\ru{Хуже}.}}) \textquotedblleft{}Cheaper\textquotedblright{}? (\emph{\textbf{\ru{Дешевле}.}})
\end{tcolorbox}
\section[\texorpdfstring{\ru{уже} \ru{два} \ru{года} … (uzhé dva góda …)}{уже два года … (uzhé dva góda …)}]{\ru{уже} \ru{два} \ru{года}; \ru{два} \ru{года} \ru{назад} — how long, and how long ago}

\label{lesson:RU-C268-uzhe-let}



\begin{quote}
Say \textquotedblleft{}already\textquotedblright{}, \textquotedblleft{}back\textquotedblright{} and \textquotedblleft{}a year\textquotedblright{}. (\emph{Uzhé. Nazád. God.}) Then \textquotedblleft{}more\textquotedblright{} and \textquotedblleft{}less\textquotedblright{}. (\emph{Bólshe. Ménshe.}) Comparing things leads to comparing times: how long, and how long ago.
\end{quote}

\begin{grammarlens}[title={Grammar Lens: how long, in the present}]
\begin{quote}
\textbf{\ru{Я} \ru{уже} \ru{два} \ru{года} \ru{учу} \ru{русский}.} — \emph{ya uzhé dva góda uchú rússkiy} — I've been learning Russian for two years.
\end{quote}

English says \textquotedblleft{}I have been learning\textquotedblright{}. Russian keeps the present, because the learning is still going on; \textbf{\ru{уже}}, \textquotedblleft{}already\textquotedblright{}, marks how far it has come. In the same way: \textbf{\ru{Я} \ru{давно} \ru{живу} \ru{здесь}.} — \emph{ya davnó zhivú zdes} — I've lived here for a long time.
\end{grammarlens}

\begin{grammarlens}[title={Grammar Lens: how long ago, in the past}]
\textbf{\ru{Назад}}, \textquotedblleft{}back\textquotedblright{}, follows a stretch of time to mean \textquotedblleft{}ago\textquotedblright{}:

\begin{quote}
\textbf{\ru{Два} \ru{года} \ru{назад} \ru{я} \ru{был} \ru{в} \ru{Москве}.} — \emph{dva góda nazád ya byl v Moskvé} — Two years ago I was in Moscow. \textbf{\ru{Час} \ru{назад} \ru{позвонила} \ru{мама}.} — \emph{chas nazád pozvoníla máma} — Mum rang an hour ago.
\end{quote}
\end{grammarlens}

\begin{grammarlens}[title={Grammar Lens: \ru{год}, \ru{года}, \ru{лет}}]
After a number, the word for a year changes: \textbf{\ru{один} \ru{год}}; \textbf{\ru{два}, \ru{три}, \ru{четыре} \ru{года}}; from \textbf{\ru{пять}} to \textbf{\ru{двадцать}}, \textbf{\ru{лет}}. Ask someone's age with it: \textbf{\ru{Сколько} \ru{тебе} \ru{лет}?} — \emph{skólko tebé let?} And answer: \textbf{\ru{Мне} \ru{двадцать} \ru{лет}.} — \emph{mne dvádtsat let} — I'm twenty.
\end{grammarlens}

\begin{cousinweb}
\textbf{\ru{лет}} is the counting form of \textbf{\ru{лето}}, \textquotedblleft{}summer\textquotedblright{}. EtymDB 2.1 traces \emph{\ru{лето}} to Proto-Slavic \emph{*lěto}, which already meant \textquotedblleft{}year\textquotedblright{}, so a long count of years is a count of summers. \textbf{\ru{год}} goes back through Old East Slavic \emph{godŭ}, \textquotedblleft{}time, year\textquotedblright{}, to Proto-Slavic \emph{*godŭ}, \textquotedblleft{}the suitable time\textquotedblright{} (EtymDB).
\end{cousinweb}

\subsection*{Your turn}
\begin{itemize}
\raggedright
  \item \emph{Say it:} \emph{Ya uzhé dva góda uchú rússkiy.}
  \item \emph{Say it:} \emph{Dva góda nazád ya byl v Moskvé.}
  \item \emph{Say it:} \emph{odín god, dva góda, pyat let}
  \item \emph{Answer:} \emph{Skólko tebé let?} \emph{Mne …}, then \emph{Skólko let ty uchísh rússkiy?}
\end{itemize}

\begin{tcolorbox}[breakable,colback=teal!4,colframe=teal!35!black,title={Before you move on}]
\textquotedblleft{}I've been learning Russian for two years\textquotedblright{}: which tense? (\textbf{The present}, with \emph{\textbf{\ru{уже}}}.) \textquotedblleft{}Two years ago\textquotedblright{}? (\emph{\textbf{\ru{Два} \ru{года} \ru{назад}.}}) \textquotedblleft{}Five years\textquotedblright{}? (\emph{\textbf{\ru{Пять} \ru{лет}.}})
\end{tcolorbox}
\section[(review)]{Review — better than, and years ago}

\label{lesson:RU-C268-review}



\begin{quote}
Say \textquotedblleft{}the train is faster than the bus\textquotedblright{}. (\emph{Póyezd bystréye, chem avtóbus.})
\end{quote}

\begin{grammarlens}[title={What you've built}]
Two lessons, and an experience can be measured. \textbf{-\ru{ее}} and \textbf{\ru{чем}} for most comparisons, and six to learn whole: \textbf{\ru{больше}}, \textbf{\ru{меньше}}, \textbf{\ru{лучше}}, \textbf{\ru{хуже}}, \textbf{\ru{дороже}}, \textbf{\ru{дешевле}}. The present with \textbf{\ru{уже}} for what is still going on, the past with \textbf{\ru{назад}} for when it happened, and \textbf{\ru{год}}, \textbf{\ru{года}}, \textbf{\ru{лет}} to count.

A trip, measured:

\begin{quote}
\textbf{\ru{Два} \ru{года} \ru{назад} \ru{я} \ru{был} \ru{в} \ru{Москве}. \ru{Поезд} \ru{был} \ru{дешевле}, \ru{чем} \ru{самолёт}, \ru{и} \ru{мне} \ru{было} \ru{интересно}. \ru{Я} \ru{уже} \ru{два} \ru{года} \ru{мечтаю} \ru{вернуться}.}
\end{quote}
\end{grammarlens}

\subsection*{Your turn}
\begin{itemize}
\raggedright
  \item \emph{Say it:} \emph{Dva góda nazád ya byl v Moskvé.}
  \item \emph{Say it:} \emph{Póyezd byl deshévle, chem samolyót.}
  \item \emph{Say it:} \emph{Ya uzhé dva góda mechtáyu vernútsya.}
  \item \emph{Answer:} \emph{Shto lúchshe, po-tvóyemu: móre íli góry?} \emph{Po-tvóyemu} is \textquotedblleft{}in your view\textquotedblright{}. Answer with \emph{po-móyemu} and one comparison
\end{itemize}

\begin{tcolorbox}[breakable,colback=teal!4,colframe=teal!35!black,title={Before you move on}]
\textquotedblleft{}More expensive\textquotedblright{}? (\emph{\textbf{\ru{Дороже}.}}) \textquotedblleft{}An hour ago\textquotedblright{}? (\emph{\textbf{\ru{Час} \ru{назад}.}}) \textquotedblleft{}Twenty years\textquotedblright{}? (\emph{\textbf{\ru{Двадцать} \ru{лет}.}})
\end{tcolorbox}
